\section{Variable twists and the complete twisted column comparison}
\label{sec:twisted-comparison}

The reference comparison TG12 extends to every member of the actual
twist family. The construction below proves the local chart statements
before assembling their global comparisons. In particular it retains
the prescribed frame depth, the inverse in the level character, and the
actual moving element in the terminal period map.

We use the supplied arithmetic host, the true column datum \(I\), its
literal \texttt{Ok} and \texttt{Pins}, and its \texttt{Later} data.
The laws may be taken from TL13. Fix a word \(u\) and a twist \(a\).
Write \(Q_E(a)\) for its selected primes of \(E\), and \(Q_F(a)\)
for their underlying primes of \(F\). The selection contains no conjugate
pair; these primes avoid the type and auxiliary primes. At every
\(P\in Q_E(a)\), the original \(\pi=\pi_u\) is a unit.
Completions are interpreted through the actual product
\(E\otimes_F F_v=\prod_{P\mid v}E_P\).
The notation \(\Delta\) always denotes the integral Weyl lattice used
in CV and CR. We never replace it by its dual Weyl lattice.

\paragraph{TX1: completing a rational level and its character.}
Let \(P\) be a finite prime, let \(H\subset\operatorname{GL}_n(O_{E,(P)})\),
and let \(\lambda:H\to K^\times\) be a character. Suppose that the
principal congruence subgroup of depth \(N\) is contained in \(H\)
and that \(\lambda\) is trivial on it. Let \(d\ne0\) have valuation at
least \(N\). In the completed integral general linear group put
\[
 K_P(d)=\{z:z-I\in dM_n(O_{E_P})\},\qquad
 \widehat H=H K_P(d).
\]
Here the rational matrices in \(H\) are embedded by the prescribed
completion map. Conjugating \(z-I\) by an integral matrix with integral
inverse preserves \(dM_n(O_{E_P})\). Thus \(K_P(d)\) is normal in
the completed integral general linear group, and \(\widehat H\)
is a subgroup.

There is a unique character
\[
 \widehat\lambda(gz)=\lambda(g)\quad(g\in H,\ z\in K_P(d)).
 \tag{TX1.1}
\]
Indeed two such expressions imply that the rational matrix
\(g^{-1}g'\) is congruent to \(I\) to depth \(N\). It therefore
belongs to the principal subgroup on which \(\lambda\) is trivial.
Normality proves multiplicativity. The same argument gives the exact
intersection
\[
 \widehat H\cap\operatorname{GL}_n(E)=H
 \quad\text{among the rational matrices integral with integral inverse.}
 \tag{TX1.2}
\]
It also shows that any exponent annihilating \(\lambda\) annihilates
\(\widehat\lambda\). The subgroup is open since it contains \(K_P(d)\).
These are the source constructions \texttt{qbU} and \texttt{qbChr};
the rational intersection and character restriction are part of the
proof, rather than assumptions on an unspecified completed level.

\paragraph{TX2: the actual terminal depth and the paired-prime diagonal.}
Let \(c=\texttt{topDep}(R,I.\texttt{cst}(a).M)\).
The terminal rational level is generated by the block extension of
the row level and the principal subgroup defined by divisibility
of the matrix discrepancy by \(q_P^{\,c+1}\).
The latter is normal in the integral general linear group, so this
generated subgroup is their product. Its character is the row
character on the first factor and is trivial on the principal factor.
Consequently the precise source choice
\[
 d_{\rm top}(a,P)=q_P^{\,c+1}
 \tag{TX2.1}
\]
satisfies the hypotheses of TX1. This uses the source valuation of
\(q_P\): in additive \(P\)-valuation its principal depth is
\(N_P=v_P(q_P^{\,c+1})\). One must not identify \(N_P\) with
\(c+1\) without a degree-one unramified hypothesis. Nor do we
replace the prescribed exponent \(c+1\) by an unspecified large
integer.

Put \(D_\pi=\operatorname{diag}(1,\ldots,1,\pi)\). At each
selected \(P\) prescribe the component \(D_\pi\); at its conjugate
prescribe its inverse conjugate transpose; prescribe identity at all
other primes. Because no selected prime is conjugate to another, these
prescriptions are compatible. The resulting finite adele \(t_a\)
is unitary. Both paired components are integral with integral inverse,
since \(\pi\) is a unit at these primes. All prescribed type and
auxiliary reductions are identity. Thus
\[
 t_a\in K^1_{ST,SG},\qquad
 \rho_P(t_a)=D_\pi\ (P\in Q_E(a)),\qquad
 (t_a)_v=1\ (v\notin Q_F(a)).
 \tag{TX2.2}
\]
This constructs the required element directly. The rational matrix
\(D_\pi\) itself is not being asserted to be unitary: unitarity is
the paired-prime condition on the entire adele.

\paragraph{TX3: the framed level dictionary, including its inverse.}
Let \(K_0\) be the reference terminal chart, with level \(U_0\),
character \(\chi_0\), and its actual standard companion \(k\).
The reference reading gives
\[
 x\in U_0\ \Longleftrightarrow\ k^{-1}xk\in\operatorname{levU}(\Lambda),
 \qquad
 \chi_0(x)=\sigma(\chi_\Lambda(h))
 \tag{TX3.1}
\]
when the auxiliary reduction of \(k^{-1}xk\) is \(h^{-1}\).
At \(Q_E(a)\), \(k\) is identity. Complete each actual terminal
frame level \((\Theta_P,\lambda_P)\) by TX1 and TX2, and set
\[
 U_a=U_0\cap\bigcap_{P\in Q_E(a)}
 \{x:\rho_P(t_a^{-1}xt_a)\in\widehat\Theta_P\},
 \quad
 \chi_a(x)=\chi_0(x)
 \left(\prod_{P\in Q_E(a)}
       \sigma\widehat\lambda_P(\rho_P(t_a^{-1}xt_a))\right)^{-1}.
 \tag{TX3.2}
\]

Suppose \(ykt_a\) and \(yxkt_a\) realize framed points \(p_0,p_1\).
Put \(z=(kt_a)^{-1}x(kt_a)\). Since \(t_a\) lies in the congruence
kernel of the reference level, it normalizes that level and has
trivial auxiliary reduction. Thus the underlying decoration and
auxiliary conditions on \(z\) are exactly the reference ones.
If \(A_P,A'_P\) are the two rational frames, equality of their
underlying local lattices gives
\[
 \rho_P(z)=A_P^{-1}A'_P
       \in\operatorname{GL}_n(O_{E,(P)}).
 \tag{TX3.3}
\]
By (TX1.2), the new level condition is precisely
\(\theta_P=(A_P^{-1}A'_P)^{-1}\in\Theta_P\).
The frame action is \(A_P\mapsto A_P\theta_P^{-1}\).
Together with the reference auxiliary action this proves
\[
 x\in U_a\quad\Longleftrightarrow\quad
 p_1=(1,h,\theta)\,p_0
 \quad\text{for some actual finite-level element }(h,\theta).
 \tag{TX3.4}
\]
The local completion character in (TX3.2) is evaluated at
\(\theta_P^{-1}\). Its outer inverse therefore gives
\(\lambda_P(\theta_P)\), and (TX3.1) gives
\[
 \chi_a(x)=\sigma\bigl(\chi_\Lambda(h)\chi_\Theta(\theta)\bigr).
 \tag{TX3.5}
\]
In particular \(\chi_a(x)=\chi_0(x)\) for the same matrix \(x\)
whose components at \(Q_F(a)\) are identity. This is a pointwise
equality of characters on a specified subgroup, not an identification
of characters on unrelated isomorphic groups.

\paragraph{TX4: constructing the framed chart, rather than postulating it.}
The level \(U_a\) is open, lies in the prescribed lattice stabilizer,
and contains a principal congruence subgroup at every relevant place.
It therefore defines the required step level with the same base
lattice \(L_\pi\). Outside \(Q_F(a)\) its local group and its lattice
are literally those of \(K_0\).

For each preferred framed-point class \(c\), DR3 supplies a full
\texttt{IsPtOf} realizer \(g_c\) of its extension to the actual \(SQ\)
point. Choose
\[
 y_c=g_c(kt_a)^{-1}.
 \tag{TX4.1}
\]
The point reading is then true for these very representatives.
The level dictionary (TX3.4) proves injectivity of the associated
map from point orbits to unitary double cosets.

Here is the surjectivity argument, including the rational-frame issue.
Start with an arbitrary unitary adele \(y\). At the finitely many
selected primes approximate the columns of \(ykt_a\), which give
completed bases of full lattices. Rational
matrices are dense in these completions, and the conditions of being
a basis, preserving each prescribed congruence, and having discrepancy
in \(K_P(d_{\rm top})\) are open. Choose rational matrices in these
neighbourhoods. Prescribe the corresponding small correction at each
\(P\), its inverse conjugate transpose at \(\bar P\), and identity
elsewhere. The paired-prime splicing of AD4 gives a correction
\(z\in K^1_{ST,SG}\) with precisely these components. The element
\(x=kt_a z(kt_a)^{-1}\) lies in \(U_a\), by (TX3.2), and
\[
 yxkt_a=ykt_a z
\]
has the chosen rational frames. The intersection construction for its
lattices and reductions gives a
valid framed point, hence one of the point classes. This proves
surjectivity.

The same dictionary gives the stabilizer identity
\[
 \operatorname{im}(\Gamma_c\longrightarrow U(E))
       =U(E)\cap y_cU_a y_c^{-1}.
 \tag{TX4.2}
\]
Injectivity follows from the actual point action: once its rational
component and the point are fixed, its auxiliary reductions and frame
components are forced. Pulling back the standard completed lattices
through \(y_c\) recovers the source lattice assigned to \(c\).
Thus the orbit bijection, the integral lattices, and the stabilizers
are all attached to (TX4.1). These are the substantive clauses of
\texttt{exists\_clsG\_top}.

\paragraph{TX5: terminal readings, roots, base points, and extension.}
At the last step the index set is a singleton and the coupling space
is empty. TX4 consequently gives the actual single chart of
\texttt{chart\_twm}. Its \texttt{Reads} clauses are as follows.
Its base lattice is \(L_\pi\); its companion is \(kt_a\); its
components outside the type and twist primes are unchanged; its
inverse has the required \(\bar\pi\) last row at type primes; and
its representatives, local move dictionary, character dictionary,
and stabilizers are exactly (TX4.1), (TX3.4), (TX3.5), and (TX4.2).
This proves \texttt{reads\_top}.

The reference character has finite exponent \(e_0>0\). Each actual
frame character has finite exponent; a common multiple \(e_Q>0\)
annihilates their product and its inverse. Thus \(e_0e_Q\)
annihilates \(\chi_a\), proving \texttt{root\_top}.
If \(gkt_a\) and \(ykt_a\) realize the same preferred point,
apply (TX3.4) with the identity finite-level element. It gives
\[
 g^{-1}y\in U_a,\qquad \chi_a(g^{-1}y)=1.
 \tag{TX5.1}
\]
In particular it gives \texttt{top\_base}, with its exact order
of multiplication.

Finally let \(g'\) realize an actual row point \(d\), and suppose
\(\operatorname{IsExt}(z_0,g',g)\), where
\(z_0=\pi e_{\rm last}\). Since \(\pi\ne0\), the extension
equations force \(g=\operatorname{diag}(g',1)\). The underlying
decoration of \(gk\) is the terminal extension of \(d\), by TG6:
\[
 N=N_{\rm row}\oplus O_E,\quad
 L=N_{\rm row}\oplus\pi O_E,\quad
 L^\sharp=N_{\rm row}\oplus\bar\pi^{-1}O_E.
 \tag{TX5.2}
\]
Its functional is the reduction of \(\bar\pi\) times the last
coordinate. The auxiliary reading uses the prescribed congruence
\(\pi\equiv1\) there. Right multiplication by \(t_a\in K^1_{ST,SG}\)
preserves all these data. At \(P\in Q_E(a)\), however, it supplies
the required point frame:
\[
 \operatorname{diag}(A_P,1)\,D_\pi
       =\operatorname{diag}(A_P,\pi).
 \tag{TX5.3}
\]
Thus \(gkt_a\) realizes \(\widetilde F(d)\), including every twist
frame. Acting matrices are \(\operatorname{diag}(\gamma,1)\);
point frames are (TX5.3). This proves the full extension clause
of \texttt{chart\_twm}, and hence its \texttt{TopBase} implication.
No separate terminal point or realization premise remains.

\paragraph{TX6: middle charts with the actual coupling stabilizers.}
For \(1\le i<m\), a point has an index \(b\), an auxiliary coupling
class \(\bar C\), and at each twist prime a coupling coset \(C_P\).
The latter is empty away from the penultimate stage. At that stage
it has the form
\[
 C_P=z_P+q_P^{M}O_{E,(P)}^n,
 \qquad z_P\in q_P^{-M}O_{E,(P)}^n,
\]
as follows by applying the actual rational frame to the two
\texttt{locLat} bounds in \texttt{ybSet}. Define a chart fibre by
fixing these data. There are finitely many fibres: the index set and
auxiliary residue spaces are finite, and the twist cosets are indexed
by the finite quotient
\(q_P^{-M}O_{E,(P)}^n/q_P^{M}O_{E,(P)}^n\).
Its additive depth is \(2M v_P(q_P)\), with the actual valuation.

In the fibre \(\beta\), cut the auxiliary level by the stabilizer
of \(\bar C\), and cut the rational twist level by the stabilizers
of the actual \(C_P\). The uncut twist level is a subgroup of
\(\operatorname{GL}_n(O_{E,(P)})\): this is part of the definition
of \texttt{LevI}, also before the penultimate stage. Its character
is trivial on a sufficiently deep principal subgroup, by the actual
\texttt{ThetaOpen} statement \texttt{lvQ\_open}. Here is the
underlying reason. Bound the denominators of the index matrix
\(h\) and \(h^{-1}\). If \(g\) is sufficiently close to identity,
then \(h^{-1}\operatorname{diag}(g,I)h\) and its inverse satisfy
the defining twist-level congruences, with identity residue matrix.
The zero extension is integral, so the additive character is trivial
on its phase; the residue row character is also \(1\).
The actual group and character readings in \texttt{Ok} then put
\(g\) in the twist level with character \(1\).

Increase this depth until
\((g-I)z_P\in q_P^{M}O_{E,(P)}^n\).
Integral \(g,g^{-1}\) already preserve the direction lattice, so
this additionally stabilizes \(C_P\); the empty coset gives no
extra restriction. Thus the cut level and its character satisfy
TX1. Character triviality here comes from the proved congruence
statement, rather than from finite order alone.

Choose an admissible companion \(k_\beta\), with identity components
at \(Q_E(a)\), as supplied in TX7. Apply DR3 to the actual preferred
points and use representatives \(y_c=g_c k_\beta^{-1}\).
The proof of TX4 now applies with these stabilizer cuts: a local
move is allowed precisely when it preserves the fixed coupling
data. The criterion is necessary by applying the move to the
point, and sufficient because the auxiliary and twist reductions
recover the coupling class. Rational approximation inside the
cut open subgroups proves double-coset surjectivity. This gives
finite charts, actual readings, lattices and stabilizers at every
middle step, with their actual finite-order characters.

The subsequent \texttt{FunKilled} normalization changes only the
value module and base functional by the FC4 construction on
\(L_\pi^\sharp/L_\pi\).
It changes neither the local groups nor their characters, the
representatives, the coupling stabilizers, or the companion
\(k_\beta\). Hence all readings persist. FC9 then transports
these same charts to the Hermitian description, giving
\texttt{SameCar}, rather than choosing unrelated Hermitian charts.

\paragraph{TX7: matching a middle chart to the reference away from \(Q\).}
We justify the admissible companions and the reference comparison
used in TX6. At an auxiliary prime, a nonzero actual index
representative lies in the identity double coset. The induction
starts with the zero-step identity, uses nonvanishing of the
preceding kernel and double-coset transport at each successor,
and ends with the prescribed final identity. The finite-sum
factorization and identity nonvanishing needed for this argument
are the local statements of CI7 and CI13. Thus, for a twisted
representative \(h_a\) and a reference representative \(h_0\),
write \(h_a=r_a\kappa_a\), \(h_0=r_0\kappa_0\) with \(r\) in
the left level and \(\kappa\) in the right level. The equality
\[
 h_a=(r_ar_0^{-1})h_0(\kappa_0^{-1}\kappa_a)
 \tag{TX7.1}
\]
conjugates the projected level groups and their characters.
It follows by substituting in the defining stabilizer equation;
the kernel eigencharacter is unchanged under the corresponding
affine change of summation variables. At the penultimate stage
both representatives are literally identity, so this conjugator
can be chosen to be identity there.

At that exceptional stage one must also match the coupling class.
Choose a vector \(y\) representing the given auxiliary coupling
residues, and form in the reference datum the coset
\[
 Y_0=y+\operatorname{locLat}_{SG}(N;0).
 \tag{TX7.2}
\]
There is no reference \(Q\)-condition. The actual reduction
formula shows that \(Y_0\) has exactly the given auxiliary coupling
vector. Move this point to its preferred representative. If the
auxiliary component of this move is \(g_0\), use \(g_0^{-1}\)
as the additional conjugator. The equality
\(\bar C_0=g_0^{-1}\bar C\) shows directly that the two coupling
stabilizers are conjugate. Before the penultimate stage the
coupling spaces are empty and this correction is unnecessary.

Lift the resulting finite auxiliary matrices to a unitary adele
\(s_\beta\), identity at type primes and outside the auxiliary
places, using paired-prime integral lifting. This is the AD4
construction; no global surjectivity of
\(\operatorname{GL}_n(O_E)\to\operatorname{GL}_n(O_E/P)\)
is needed. Put
\[
 k_\beta=k_{0,g(\beta)}s_\beta.
 \tag{TX7.3}
\]
The conjugators preserve the standard lattices, so this is an
admissible companion and is identity at \(Q_E(a)\).
For an adele whose \(Q_F(a)\)-components are identity, its frame
move is identity. Equations (TX7.1)--(TX7.3) and the coupling
calculation then give the same membership condition and the same
character value in the chart and its reference chart. Taking
adeles supported at a single place proves the literal equality
\[
 U_{\beta,v}=U_{0,g(\beta),v}\quad(v\notin Q_F(a)),
 \qquad
 \chi_\beta(x)=\chi_{0,g(\beta)}(x)
       \quad(x_v=1\text{ for }v\in Q_F(a)).
 \tag{TX7.4}
\]
The base lattices are also the same \(L_\pi\). The reference
twist itself uses its reference charts unchanged.

\paragraph{TX8: a character on the common penultimate local group.}
A character on each chart need not extend to the group generated
by all charts. The following argument supplies that missing step.
Fix \(v\in Q_F(a)\) and the penultimate index, whose representative
is identity. For each selected \(P\mid v\), use the positive
integer \(D(P)\) in the actual local constants. Let
\(\mathcal N_v\) consist of local matrices which have a rational
approximation \(g_P\) to this depth satisfying the actual rational
middle-level relation, with character value \(z_P\).
This is exactly the source \texttt{NearSet}; approximation and
congruence are measured through its prescribed completion maps.

Two witnesses \(g_P,g'_P\) for the same matrix differ by an
integral matrix congruent to identity to depth \(D(P)\).
Both matrices and their inverses are integral: at the penultimate
stage the projected rational level lies in
\(\operatorname{GL}_n(O_{E,(P)})\). At auxiliary primes its
character is the prescribed residue character, and at twist
primes it is the actual \(\lambda_P\) on the rational level.
At this penultimate step the tail is empty, so the additive
\texttt{n1Arg} phase in \texttt{chiInv} is absent. The remaining
\texttt{lamRow} character factors through the actual residue
matrices. The positive prescribed exponent \(D(P)\) makes a
matrix congruent to identity modulo \(q_P^{D(P)}\) act trivially
in these residue matrices. Both characters are therefore trivial
on the stipulated congruence subgroup, and
\[
 z_P=z'_P.
 \tag{TX8.1}
\]
This proves independence of every rational-lift choice.

Products and inverses of witnesses are again witnesses. For
example the discrepancy identity
\[
 (AB)^{-1}(CD)
   =(B^{-1}D)\bigl(D^{-1}(A^{-1}C)D\bigr)
 \tag{TX8.2}
\]
preserves the same congruence depth because \(D,D^{-1}\) are
integral; positivity of the depth ensures the required integral
neighbourhood. The rational middle-level character is
multiplicative. Thus \(\mathcal N_v\) is a group, and
\[
 \operatorname{val}_v(x)
       =\sigma\!\left(\prod_{P\mid v}z_P\right),
 \qquad
 \psi_v(x)=\operatorname{val}_v(x)^{-1}
 \tag{TX8.3}
\]
are well-defined characters. Factors at primes not over \(v\)
are \(1\), as in the literal source definition.

Each penultimate chart group at \(v\) lies in \(\mathcal N_v\).
Indeed approximate its completed frames by rational frames while
preserving the finitely many open level and coupling conditions.
The actual local-move dictionary then supplies a rational
middle-level witness; \(k_\beta=1\) at \(P\) identifies its
congruence with the stipulated \(D(P)\)-congruence. Its character
is the inverse product of these rational values, by the same
inverse-reduction calculation as TX3. Consequently
\[
 \chi_\beta(\operatorname{emb}_v x)=\psi_v(x)
       \quad(x\in U_{\beta,v}).
 \tag{TX8.4}
\]
Different close witnesses give the same value by (TX8.1).
Let \(W_v\) be the subgroup generated by all \(U_{\beta,v}\).
Then \(W_v\subset\mathcal N_v\), and the restriction of \(\psi_v\)
is the required common character on \(W_v\).

The penultimate chart family is nonempty. Choose the index supplied
by \texttt{last\_uniq}, the base decoration, and the zero-centred
actual \texttt{locLat} coupling coset. It is allowed since
\(0\) belongs to every lattice involved. This gives a point class
and hence a chart. Its principal congruence subgroup lies in
\(W_v\), proving that \(W_v\) is open.

Each chart group has finite index in \(W_v\). All groups lie
in the same local lattice stabilizer, and each chart contains a
principal congruence subgroup. The latter is the kernel of the
action on a finite module \(L/P^rL\); its containing group
therefore has finite index. This argument establishes both the
common character and the finiteness needed below.

\paragraph{TX9: the common penultimate level and its finite coupling carrier.}
Let \(H_0\) be the reference common penultimate level of FC12 and
TG2, with its identity auxiliary frames. Define \(J\) by replacing
its local group at \(v\in Q_F(a)\) by \(W_v\), keeping all other
local groups and the base lattice unchanged. The product
decomposition of the finite adele ring gives a homomorphism
\(\operatorname{off}_Q:J\to H_0\) which replaces the
\(Q_F(a)\)-components by identity. Put
\[
 \chi_J(\kappa)=
       \chi_{H_0}(\operatorname{off}_Q\kappa)
       \prod_{v\in Q_F(a)}\psi_v(\kappa_v).
 \tag{TX9.1}
\]
Decompose a chart element into its part supported on \(Q_F(a)\)
and its part identity there. TX8 identifies the first character;
TX7 identifies the second. Thus \(\chi_J\) restricts to the
actual character of every penultimate chart. It equals the
reference common character on the same elements identity at \(Q\).

For completeness, the finite carrier can be constructed explicitly.
Let \(\rho:J\to G_{\rm aux}\) be the actual auxiliary reduction
map, and let \(w_\beta\) be the auxiliary coupling vector of
chart \(\beta\). If \(K_{\beta,v}=U_{\beta,v}\subset W_v\),
the full chart stabilizer in \(J\) is
\[
 H_\beta=\{\kappa:
       \rho(\kappa)w_\beta=w_\beta,\quad
       \kappa_v\in K_{\beta,v}\ (v\in Q_F(a))\}.
 \tag{TX9.2}
\]
Take the disjoint union of the \(J\)-orbits of
\((w_\beta,(K_{\beta,v})_v)\). Record the chart label in the
second coordinate, so distinct chart orbits remain distinct.
Equivalently this carrier is \(\coprod_\beta J/H_\beta\), with
the first map \(y_G\) recording the translated auxiliary vector
and the second map \(y_Q\) recording the labelled local cosets.
Equation (TX9.2) proves that the pair \((y_G,y_Q)\) is injective.
It also proves the required stabilizers and equivariance:
\(y_G(\kappa y)=\rho(\kappa)y_G(y)\), and \(y_Q\) is unchanged
when all \(\kappa_v\), \(v\in Q_F(a)\), are identity.
The carrier is finite by TX8 and finiteness of the auxiliary
residue spaces.

At auxiliary primes \(J\) is the exact symmetric group
\(H_{\rm Sym}\) of TG2, with its actual reduction map and
identity frame. Away from the auxiliary and twist primes it
is the reference local lattice group. This proves all the
\texttt{TwLevJ} and \texttt{SymmetricTwistedJointForms} clauses,
including the common character; no unsupported extension of
characters to a generated group is used.

\paragraph{TX10: constructing \texttt{Law0} with the integral lattice.}
At step zero the auxiliary level is independent of the twist,
every twist-frame level is the whole integral general linear
group, the coupling coset is empty, and the index is unique.
Forgetting the twist frames therefore gives a bijection
\[
 e_a:\mathrm{Cl}(u,\mu_0,0)\longrightarrow\mathrm{Cl}(u,a,0).
 \tag{TX10.1}
\]
Surjectivity follows by choosing rational bases of the localized
underlying lattice at the finitely many frame primes. If two
underlying decorations differ by a rational unitary element,
their frame discrepancy is integral with integral inverse and
hence belongs to the full zero-step frame level. This proves
injectivity on point orbits.

Preferred representatives require a further check. Their
underlying decoration is the canonical representative for the
common action on decorations. Taking the canonical representative
again has no effect. Thus the preferred decorations corresponding
under (TX10.1) are literally equal, and in particular so are
their lattices \(L^\sharp\), functionals, and auxiliary maps.
The zero-step stabilizers are trivial by the actual free-action
law proved in TL. Consequently coefficient functions on the
finite class set satisfy the equivariance condition automatically.
Pullback along \(e_a^{-1}\) gives the \(L\)-linear equivalence
\[
 E_a f(c)=f(e_a^{-1}c)
 \tag{TX10.2}
\]
between the actual rational form spaces.

Membership in the global integral lattice is pointwise
membership in the same integral \(\Delta\)-lattice attached to
the preferred \(L^\sharp\). Its equality at corresponding
classes proves, in both directions,
\[
 f\in\mathcal L_{u,\mu_0,0}
       \quad\Longleftrightarrow\quad
 E_af\in\mathcal L_{u,a,0}.
 \tag{TX10.3}
\]
The coefficient embeddings commute with the same class
reindexing, proving the literal embedding clause of
\texttt{Law0}. At \(a=\mu_0\), the canonical class equivalence
is identity; injectivity of the form embedding forces the
chosen comparison to be the specified \(e_{\mu_0}\).
More generally, the embedding clause and injectivity of the
form embedding force every permissible \(E_a\) to equal
(TX10.2). The conclusions therefore hold for any supplied
\(\ell_0:\texttt{Law0}\), not only for the witness just
constructed.
Thus \texttt{Law0} is derived from the given \texttt{Ok} and laws.
It is not an additional arithmetic or integral-comparison premise.

\paragraph{TX11: all-twist chains, values and the same reference object.}
Use TG12 for the reference columns and retain that object throughout.
For each \((u,a)\), choose the middle charts of TX6--TX7, the
terminal chart of TX5, and the common penultimate level of
TX8--TX9. Make the terminal choice first and use that same choice
when collecting the steps. At \(a=\mu_0\) use the original
reference charts. At step zero use (TX10.1); outside the operative
range the chart index can be a singleton, since its reading
conditions have no instances.

Set the class lattice equal to the actual \(I.\texttt{latAt}\),
and use the identity map on frame coordinates. The coordinate
maps to the reference Hermitian spaces have the same underlying
frames by \texttt{StdF}; thus \texttt{IdY} holds literally.
The \texttt{TwChTied} data now use these chart sets, these
lattices, the actual point stabilizers, and the finite carrier
of TX9. Its level containment and stabilizer clauses are
(TX9.2). The reference comparison map \(g(\beta)\), the local
equalities (TX7.4), and prime disjointness give
\texttt{TwChJointS}. TX9 gives its symmetric joint forms.

The type-prime comparison can be chosen with both the prime
identification and residue-field identification equal to identity.
For the diamond operator at step zero, the two moving elements
act between equal underlying preferred decorations by TX10.
Their quotient fixes that decoration. Add the uniquely determined
frame correction in the full zero-step frame level; freeness
forces both their rational and auxiliary components to agree.
The zero-step frame character is trivial, so their full
transport factors also agree. This proves the diamond comparison.
The zero-step complex form spaces are related by (TX10.2),
since freeness removes the stabilizer condition. All other
type-prime parameters are the actual unchanged parameters.
This constructs \texttt{TwAtT}, with the stipulated residue
identity, rather than assuming it.

Apply CV9 to each actual reading and \texttt{IdY}. This supplies
the complex image, rational image, and the place-by-place
integral \(\Delta\)-lattice assertion. Restricting the latter
to the required good coefficient places gives exactly the
\texttt{ValuationCompatibleChart} field of \texttt{TwFormsTiedS}.
FC18 applies to the normalized charts because TX6 retained
\texttt{FunKilled}, all readings, and the same type-prime map.
It supplies the diamond clause. FC7 applies to these same frame
functions and actual integral form space, supplying its metric
and norm clauses. Equations (TX7.4) and (TX9.1) supply every
off-\(Q\) character clause.

The reference frame equivalence is identity, and TX10 identifies
the chosen zero-step equivalence with \(e_{\mu_0}\). The standard
data therefore have exactly the source \texttt{StdW} fields:
the actual \(e_{\rm Cl,0}\), identity frame and prime maps,
\(I.\texttt{latAt}\), the prescribed residues, all readings,
the terminal base and extension clauses of TX5, and
\texttt{FunKilled}. This proves \texttt{tw\_chains\_std} and
\texttt{tw\_ties\_with\_points}, with \texttt{Law0} supplied by
TX10 and no separately assumed chain or base-point object.

\paragraph{TX12: the actual level maps and zero/reference compatibility.}
On each class use the quotient projections of FC5, with their
actual kernels, for the level and kernel-level modules.
The type-prime residue map is identity under TX11's
identification. Thus FC6 identifies the transported functional.
NB7--NB8 apply to each actual up/down move: DR3 supplies its
initial full realizer; the supported splice preserves every
other type functional, auxiliary reduction, twist frame and
coupling label. The original moving element therefore gives
the required transport factor. This proves the up and down
naturality clauses for these very projections. The integral
basis clause is FC17 applied to CV9. Hence the resulting level
data are \texttt{SymmetricTwistedLevelData}, and their
\texttt{StdLv} projections are the literal quotient maps.

At step zero TX10 identifies the preferred decorations and
all their local kernels. Raw up lattices depend only on the
residue class: if \(x-x'\in PN\), then
\(P^{-1}(x-x')\subset N\), so
\(N+P^{-1}x=N+P^{-1}x'\). This also identifies the invalid
branches of the total move definition. The down kernel
construction and the diamond rescaling likewise depend only
on the same underlying decoration. Their preferred
representatives agree by canonical normalization. As in TX11,
freeness equates the rational and auxiliary components of the
moving elements; their frame characters are trivial. Thus
the class maps and the transport maps agree, not merely the
unframed neighbours.

The \(B_w\) assertion at step zero is also integral: (TX10.2)
preserves the pointwise \(\Delta\)-lattice of each equal
\(L^\sharp_{\rm ker}\), and localization preserves this
equivalence in both directions. At the reference twist all
frame maps are identity. The residue-field map is identity
because it agrees on the images of the surjective quotient
\(O_E\to O_E/P\). Every quotient projection, functional,
move, transport, diamond operator and \(B_w\)-condition then
agrees. These arguments prove the full \texttt{TwLvSameS},
including both its step-zero and reference-twist parts.

\paragraph{TX13: the actual twisted terminal row chart.}
The row point datum has unique index and empty coupling.
Its base lattice is \(O_E^{\,n}\), its type-prime set is empty,
and its twist-frame level is the inverse image of the actual
finite row level under reduction at \(P\). It therefore
contains the depth-one principal subgroup, on which its
character is trivial. Apply TX1 with the exact choices
\[
 k=1,\qquad t=1,\qquad d_{\rm row}(P)=q_P.
 \tag{TX13.1}
\]
Thus the row scalar uses the first power of \(q_P\), whereas
the full terminal scalar in TX2 uses power \(c+1\). A
nonuniformizer \(q_P\) may have a larger additive valuation;
its positive valuation still places its congruence subgroup
inside the depth-one reduction kernel.

DR3 realizes each actual preferred row point. Use those
realizers themselves as \(y'_c\); no companion factor occurs
in (TX13.1). Equality of rational local frames gives the
integral discrepancy used in (TX3.3), and the completed level
intersects the rational group in exactly the row frame level.
The same argument as TX4 proves the orbit bijection,
stabilizer identification and lattice pullback. It proves
\texttt{ReadsP} for the resulting \texttt{ClassesAtPrimes}.
The character is the actual auxiliary character times the
inverse product of completed row characters, so it has a
finite exponent. The base lattice agrees with the reference
row base lattice. Outside \(Q_F(a)\) the local groups agree
with the reference row groups; on the same adele identity at
\(Q_F(a)\), the characters agree. This gives the stronger
off-\(Q\) equality, and hence the required auxiliary-place
\texttt{TwPAtGS} equality.

\paragraph{TX14: row values, norms and the conjugate period map.}
The actual row weight has the prescribed common determinant
shift. Since \(t\) is odd, write \(t=2s+1\); then
\[
 m(t+1)/2=m(s+1),\qquad (t+1)/2=s+1.
 \tag{TX14.1}
\]
This identifies the two source shift conventions entry by
entry. Adding the same integer to all weight entries does
not change the polynomial natural weight, the integral
\(\Delta\)-lattice, or its Fock pairing; the determinant
factor is retained in the frame-function map.

Apply CV10 to the actual reading of TX13, with that shift.
The frame map \(e_W\) is identity on underlying frame
functions. Polynomial equivariance plus the chart dictionary
proves its complex image. The rational image is characterized
by values in the prescribed coefficient field in rational
frames. For the integral assertion, choose finitely many
class representatives and use the finite rational-coordinate
description of \(\Delta\). Clearing denominators gives the
global lattice intersection, and localizing that intersection
gives the actual local \(\Delta\)-lattice. The multipliers
are chosen outside the fixed coefficient prime at which one
is localizing: its local membership already gives the required
bound there, while finite prime avoidance, or the Chinese
remainder theorem, supplies an element with the required
orders at the other finitely many bad primes and nonzero
residue at the fixed prime. Taking this multiplier integral
also preserves every already-good place. This proves the
localization inclusion without introducing another bad place.
This is the same two-inclusion argument as CV7--CV10, now
with TX13's row representatives. It proves the literal
\texttt{ValuationAtPrimes} assertion at all places specified
by \texttt{okPT}.

In orthonormal coordinates its pairing is
\[
 \langle z,z'\rangle_P
 =\sum_c |\Gamma_c|^{-1}
     \operatorname{nfAll}_{\mathrm{sg}}(Q_c,P_c),
 \tag{TX14.2}
\]
where \(P_c,Q_c\) are the coefficient polynomials of \(z,z'\).
The order \(Q_c,P_c\) is the source Hermitian convention.
Equation (TX14.1) identifies the shifted pairing exactly.
FC7's rational-basis argument in the actual conjugation-stable
coefficient field therefore proves \texttt{NormPTieG}.
It does not interchange embeddings on arbitrary complex
coefficients.

Finally use coefficient conjugation on the actual polynomial
model to define the semilinear map \(j_M\) of FC8. Polynomial
evaluation is injective; conjugating twice gives the inverse
on the actual row-form image. Conjugating the equivariance
equations gives precisely the row automorphy equations, so
\(j_M\) maps onto that image and only that image. Expanding
the finite class sum with the same weights as (TX14.2) gives
\[
 \operatorname{period}(w,F)=\langle j_MF,w\rangle_P.
 \tag{TX14.3}
\]
This constructs the row value, norm and period clauses for
every twist; it does not replace the row chart by the
reference chart.

\paragraph{TX15: the terminal vector, moving element and evaluation for all twists.}
Fix a row class \(c\). Let \(m_c\) be the actual moving
element taking the preferred representative of its terminal
class \(\widetilde c\) to \(\widetilde F(\operatorname{rep}c)\),
and let \(\gamma_c\) be its rational unitary component.
Take \(z_0=\pi e_{\rm last}\). Then
\(\langle z_0,z_0\rangle=p\), and it spans the prescribed
orthogonal line. The raw lattices are (TX5.2), so the action
of this same moving element gives
\[
 \gamma_c L_{\widetilde c}
      =\iota(N_c)+O_Ez_0.
 \tag{TX15.1}
\]
For \(x\in L_{\widetilde c}^\sharp\) and
\(r=\langle z_0,\gamma_cx\rangle\in O_E\), the functional is
\[
 f_v(\operatorname{pr}_v x)=\operatorname{res}_v(r).
 \tag{TX15.2}
\]
To verify it on the full dual, put \(s=\bar\pi\).
Then \(sx\) lies in the intermediate self-dual lattice.
The raw last-coordinate formula and the projection reading
give \(f_v(\operatorname{pr}_v(sx))=\operatorname{res}_v(sr)\).
Both sides are \(\operatorname{res}_v(s)\) times the
corresponding side of (TX15.2). At the selected type prime
\(s\) is a unit, so cancellation proves (TX15.2).
This is a full integral-domain argument; no equality is
inferred from equality after tensoring with a field.

Let \(y_M\) and \(y'_c\) be the already chosen full and row
representatives, and let \(k\) denote the actual terminal
companion \(kt_a\) of TX5. Put
\[
 x_c=y_M^{-1}\operatorname{rat}(\gamma_c)^{-1}
                 \operatorname{ext}(y'_c),\qquad
 \kappa_c=x_c^{-1}.
 \tag{TX15.3}
\]
TX5 shows that \(\operatorname{ext}(y'_c)k\) realizes
\(\widetilde F(\operatorname{rep}c)\). Pulling back by
\(\gamma_c^{-1}\) leaves a move with rational component
identity and with the auxiliary/frame components of \(m_c\).
The terminal coupling is empty, so TX3 applies and gives
\[
 x_c\in U_M,\qquad
 \chi_M(\kappa_c)^{-1}=\sigma(\chi(m_c)),\qquad
 \operatorname{rat}(\gamma_c)y_M
         =\operatorname{ext}(y'_c)\kappa_c.
 \tag{TX15.4}
\]
This fixes both inverse orientations. Any adelic matrix
extending \(y'_c\) and fixing \(z_0\) is this block extension,
by \(\pi\ne0\); hence the last equality gives the universal
\texttt{til\_ad} clause, not only one chosen instance.

A row stabilizer \(s\) maps to
\(m_c^{-1}\operatorname{ext}(s)m_c\). Conjugating its rational
matrix by \(\gamma_c\) cancels the two moving elements and
leaves \(\operatorname{diag}(\gamma_s,1)\). This proves the
literal \texttt{IsExt} clause for \texttt{stabMap}.
Finally the definition of the column transport evaluates
at the inverse-translated frame and multiplies by
\(\sigma(\chi(m_c))\). With (TX15.4), for every frame
function \(\phi\), and every related pair of frames, it gives
\[
 e_{W,z}(\operatorname{ev}_z\phi)(c,s)
       =\chi_M(\kappa_c)^{-1}\phi(\widetilde c,s'),
 \quad
 \operatorname{IsTr}(\gamma_c)s'=s.
 \tag{TX15.5}
\]
Here \(e_{W,z}\) is identity on the underlying frame functions.
Equation (TX15.5) holds for arbitrary \(\phi\), not just for
an invariant form. Equations (TX15.1)--(TX15.5), together
with TX14, prove the full \texttt{til\_tw},
\texttt{TwPTiedS}, \texttt{TwPEvzS}, and \texttt{StdP}
conclusions for the same charts and moving elements.

\paragraph{TX16: the complete standard all-twist comparison and its exact scope.}
Let \(C\) be the reference \texttt{ColsTiedSym} constructed
in TG12. Its carriers give \texttt{CarH} on its same
Hermitian chains, and its standard level data give the
identity penultimate frames required in TX9. TX10 supplies
\(\ell_0:\texttt{Law0}\); alternatively start with any
supplied \(\ell_0\), whose chosen maps are the same by the
uniqueness proved in TX10. Apply TX11 to these objects,
TX12 to its resulting charts, and TX13--TX15 to its
resulting full terminal companions. Thus every record in
the final join has already been constructed on the same
class sets, lattices, frame maps, residue maps and characters.

To state the literal source endpoint, retain a supplied
\(\texttt{nch}:\texttt{NChoice}(D)\) and
\(\texttt{tp}:\texttt{TupTied}(D,I.\texttt{cf},\texttt{nch})\)
with the literal equality \(\texttt{tp.xi}=I.\texttt{xi}\).
The actual character reading and selected-prime law give
\[
 v\in Q_F(a)
 \ \Longleftrightarrow\
 \exists P\mid v,\ \xi_a\text{ is ramified at }P
 \ \Longleftrightarrow\
 \texttt{tp.InQ'}(a,v).
 \tag{TX16.1}
\]
The forward implication uses the nontrivial local character
at a selected twist prime; the reverse uses the selection
of the prime above every underlying ramification place.
This identifies the prime predicate itself, rather than
merely comparing cardinalities of the two exceptional sets.

Use TX11 for \(\texttt{ch},\texttt{chJ},\texttt{chM},
\texttt{forms}\), TX12 for \(\texttt{lv},\texttt{same}\),
TX9 for \(\texttt{lvJ}\), TX13--TX15 for
\(\texttt{pt},\texttt{ptG},\texttt{ptE}\).
These are exactly the ten fields of
\(\texttt{SymmetricTwistedColumnsAtPrimes}\).
Their standard clauses are TX11's \texttt{StdW},
TX12's \texttt{StdLv}, and TX15's \texttt{StdP}.
This proves
\[
 \exists\,t:
 \texttt{SymmetricTwistedColumnsAtPrimes}
       (\texttt{tp},C,I.\texttt{tw}(\ell,\ell_0,Z)),
 \qquad I.\texttt{StdTw}(\ell,Z,\ell_0,t),
 \tag{TX16.2}
\]
the actual \texttt{colsTwPTiedS\_std} conclusion.

The geometry and comparison through TX15 need neither
\texttt{NChoice} nor \texttt{TupTied}. Their appearance in
(TX16.2) is the exact original arithmetic interface:
\texttt{NChoice} includes finite discrete-valuation-ring
extensions, valuation embeddings, and basis and integral
lattice comparisons for the normalized \(M_t\)-spaces;
\texttt{TupTied} also prescribes the distinguished twists
and the characters at those same selected prime sets.
TL15's finite-idele character readings alone do not supply
these stronger arithmetic data. They remain explicit,
as do the supplied arithmetic host, \texttt{Ok},
\texttt{Pins}, and \texttt{Later}. Conversely no independent
\texttt{Law0}, realizability, chart, common-character,
twisted-chain, terminal-vector or period-comparison premise
has been introduced. The stronger
\texttt{OneLevelSymmetricColumns} assertion is not inferred
from (TX16.2); it has its own additional level restrictions.
